\begin{textsample}{POS Dim 1 – vem\_ed\_03 – Score 3.00 – vem\_ed\_03/t011.txt}  \label{ex:f1_pos_160}
INTERCÂMBIOS VIRTUAIS EM DEBATE Considerada a principal conferência sobre \textbf{intercâmbios} virtuais no mundo, a International Virtual Exchange Conference (IVEC 2020) aconteceu entre 14 e 16 de setembro.
Sediada virtualmente pela Universidade de Newcastle (Reino Unido), \textbf{contou} com 76 apresentações, 10 pôsteres, 16 workshops, 11 simpósios, um painel de discussões e duas keynote speakers: Mirjam Hauck (Open University, Reino Unido) e Thérèse Laferrière (Université Laval, Canadá).
Com mais de 500 participantes e 325 palestrantes de 30 diferentes países, a conferência mesclou sessões gravadas e ao vivo, além de \textbf{encontros} online para socialização e networking.
Osvaldo Succi Jr., coordenador dos Projetos Colaborativos Internacionais (PCI – Cesu / Coordenação de Línguas), fez \textbf{parte} do comitê de revisores de trabalhos, mediou a sessão"Is there a silver bullet to get instructors to embrace virtual exchange?"e apresentou dois trabalhos em 15 de setembro.
Uma sessão pré-gravada com o tema"Criação, crescimento e sustentabilidade de redes de \textbf{intercâmbios} virtuais"teve perguntas e respostas ao vivo para os inscritos.
Entre os palestrantes, além de Succi, estavam Maria Leonor Alves Maia (UFPE), Ana Cristina Biondo Salomão e José Celso Freire Junior (Unesp).
Os autores relataram a criação, em 2018, do plano estratégico de internacionalização em \textbf{casa} intitulado BRaVE (Brazilian Virtual Exchange Program) pela Associação Brasileira para Educação Internacional (FAUBAI).
Outra \textbf{produção} de Succi Jr. \textbf{foi} um simpósio sobre linguística aplicada e o ensino de inglês como segundo \textbf{idioma} como ferramenta para potencializar as colaborações em \textbf{intercâmbios} virtuais.
Os demais apresentadores dessa sessão \textbf{foram} Ana Cristina Biondo Salomão (Unesp), Cara Tuzzolino-Werben (Nassau Community College, EUA) e Tiffany MacQuarrie (Penn State University, EUA).
Resumos dos trabalhos e demais informações encontram-se em https://iveconference.org/.
Internacionalização tropicalizada Além da participação no IVEC, Succi Jr. fez apresentação no Congresso Internacional de Educação e Tecnologias: Encontro de Pesquisadores em Educação a Distância (CIET: EnPED).
O trabalho"Internacionalização em \textbf{casa} por meio de Intercâmbios Virtuais: tropicalizando as experiências para o Brasil"\textbf{foi} discutido no \textbf{ambiente} virtual de aprendizagem Moodle, em 28 de agosto.
A apresentação em português \textbf{foi} legendada em Libras, espanhol e inglês e está disponível em: https://bit.ly/2Sa7FgP

% matched lemmas: ambiente, casa, contar, encontro, idioma, intercâmbio, parte, produção, ser
\end{textsample}
